% Options for packages loaded elsewhere
\PassOptionsToPackage{unicode}{hyperref}
\PassOptionsToPackage{hyphens}{url}
\documentclass[
  ignorenonframetext,
]{beamer}
\newif\ifbibliography
\usepackage{pgfpages}
\setbeamertemplate{caption}[numbered]
\setbeamertemplate{caption label separator}{: }
\setbeamercolor{caption name}{fg=normal text.fg}
\beamertemplatenavigationsymbolsempty
% remove section numbering
\setbeamertemplate{part page}{
  \centering
  \begin{beamercolorbox}[sep=16pt,center]{part title}
    \usebeamerfont{part title}\insertpart\par
  \end{beamercolorbox}
}
\setbeamertemplate{section page}{
  \centering
  \begin{beamercolorbox}[sep=12pt,center]{section title}
    \usebeamerfont{section title}\insertsection\par
  \end{beamercolorbox}
}
\setbeamertemplate{subsection page}{
  \centering
  \begin{beamercolorbox}[sep=8pt,center]{subsection title}
    \usebeamerfont{subsection title}\insertsubsection\par
  \end{beamercolorbox}
}
% Prevent slide breaks in the middle of a paragraph
\widowpenalties 1 10000
\raggedbottom
\AtBeginPart{
  \frame{\partpage}
}
\AtBeginSection{
  \ifbibliography
  \else
    \frame{\sectionpage}
  \fi
}
\AtBeginSubsection{
  \frame{\subsectionpage}
}
\usepackage{iftex}
\ifPDFTeX
  \usepackage[T1]{fontenc}
  \usepackage[utf8]{inputenc}
  \usepackage{textcomp} % provide euro and other symbols
\else % if luatex or xetex
  \usepackage{unicode-math} % this also loads fontspec
  \defaultfontfeatures{Scale=MatchLowercase}
  \defaultfontfeatures[\rmfamily]{Ligatures=TeX,Scale=1}
\fi
\usepackage{lmodern}
\ifPDFTeX\else
  % xetex/luatex font selection
\fi
% Use upquote if available, for straight quotes in verbatim environments
\IfFileExists{upquote.sty}{\usepackage{upquote}}{}
\IfFileExists{microtype.sty}{% use microtype if available
  \usepackage[]{microtype}
  \UseMicrotypeSet[protrusion]{basicmath} % disable protrusion for tt fonts
}{}
\makeatletter
\@ifundefined{KOMAClassName}{% if non-KOMA class
  \IfFileExists{parskip.sty}{%
    \usepackage{parskip}
  }{% else
    \setlength{\parindent}{0pt}
    \setlength{\parskip}{6pt plus 2pt minus 1pt}}
}{% if KOMA class
  \KOMAoptions{parskip=half}}
\makeatother
\setlength{\emergencystretch}{3em} % prevent overfull lines
\providecommand{\tightlist}{%
  \setlength{\itemsep}{0pt}\setlength{\parskip}{0pt}}
% Slide layout defaults for warning-clean scientific decks.
\usepackage{etoolbox}
\IfFileExists{xurl.sty}{\usepackage{xurl}}{}
\IfFileExists{seqsplit.sty}{\usepackage{seqsplit}}{\newcommand{\seqsplit}[1]{#1}}
\protected\def\breakseq#1{\seqsplit{#1}}
\protected\def\breaktt#1{\begingroup\ttfamily\seqsplit{#1}\endgroup}
\setlength{\emergencystretch}{6em}
\tolerance=5000
\hbadness=10000
\hfuzz=1pt
\setlength{\tabcolsep}{2pt}
\AtBeginEnvironment{longtable}{\tiny\renewcommand{\arraystretch}{0.86}\setlength{\tabcolsep}{1pt}}
\AtBeginEnvironment{tabular}{\tiny\renewcommand{\arraystretch}{0.86}\setlength{\tabcolsep}{1pt}}

% Natbib and cross-reference fallbacks — slides don't load natbib
% or cleveref, but combined-PDF manuscript prose may emit these
% commands. The fallback renders citations as a bracketed key list
% and cross-references as detokenized labels so slides don't fail on
% undefined control sequences or raw underscores. \providecommand is
% a no-op if packages load later.
\providecommand{\citep}[1]{[#1]}
\providecommand{\citet}[1]{#1}
\providecommand{\citealp}[1]{#1}
\providecommand{\citeauthor}[1]{#1}
\providecommand{\citeyear}[1]{#1}
\providecommand{\cref}[1]{\texttt{\detokenize{#1}}}
\providecommand{\Cref}[1]{\texttt{\detokenize{#1}}}

\newtheorem{proposition}{Proposition}
\newtheorem{hypothesis}{Hypothesis}
\usepackage{bookmark}
\IfFileExists{xurl.sty}{\usepackage{xurl}}{} % add URL line breaks if available
\urlstyle{same}
% fallback for those not using the hyperref driver hyperxmp:
\makeatletter
\@ifundefined{xmpquote}{\newcommand{\xmpquote}[1]{#1}}{}
\makeatother
\hypersetup{
  hidelinks,
  pdfcreator={LaTeX via pandoc}}

\author{\texorpdfstring{}{}}
\date{}

\begin{document}

\section{\texorpdfstring{Conclusion: Evidence Landscape, Methodological
Limitations, and Research Agenda
\label{sec:conclusion}}{Conclusion: Evidence Landscape, Methodological Limitations, and Research Agenda }}\label{conclusion-evidence-landscape-methodological-limitations-and-research-agenda}

\begin{frame}[allowframebreaks]{Summary}
\protect\phantomsection\label{summary}
This work demonstrates a first-generation prototype infrastructure for
citation-weighted evidence mapping and hypothesis triage across a
rapidly growing scientific field. By combining multi-source retrieval
(\(N = 817\) papers from three databases, inclusion from 2000 onward),
LLM-based assertion extraction encoded as provenance-bearing
nanopublications, and citation-weighted triage scoring, we produce a
queryable, RDF-compatible knowledge graph that maps the evolving
evidence landscape for eight core Active Inference claims. The system
demonstrates the feasibility of automated living reviews, while clearly
delineating the boundaries of current model capabilities.

All assertions and hypothesis scores are machine-generated; outputs
always report evidence mapping and triage, not scientific
confirmation---there is no live gate that suppresses or changes output
based on the validation metrics below, which are reported as
reproducibility context, not an enforced pass/fail threshold. A
stratified rule-based reference-annotator agreement study (\(n = 200\))
yields inter-rule \(\kappa = 0.554\) and pipeline precision 0.000,
recall 0.000 against a deterministic keyword-rule reference; direction
agreement between that reference and the pipeline is only
\(\kappa = -0.055\). These are reproducibility signals, not accuracy
against human labels---the pre-declared human gold-standard baseline
below has not yet been collected.
\end{frame}

\begin{frame}[allowframebreaks]{Constraints and Methodological Scope}
\protect\phantomsection\label{constraints-and-methodological-scope}
Several conscious design constraints scope these findings.

\begin{block}{Keyword Classifier Resolution}
\protect\phantomsection\label{keyword-classifier-resolution}
The keyword-based classifier operates over 200+ keyword indicators
distributed across 8 domain categories (74 mathematical indicators in A1
alone), using a deterministic priority system that routes papers to
specific application domains (C1--C5) before testing tools (B), formal
theory (A1), and the qualitative philosophy catch-all (A2).
Word-boundary-aware matching reduces partial-match false positives, but
keyword-based methods cannot capture semantic nuance: papers using novel
terminology or discussing cross-domain topics without standard
vocabulary risk misclassification. Residual A2 concentration should be
interpreted as a ceiling on broad theoretical generality rather than a
literal measure of philosophical focus. An embedding-based classifier
trained on a labeled subset would provide a quantitative upper bound on
the fraction of A2 papers that merit redistribution.
\end{block}

\begin{block}{Citation Network Coverage Gaps}
\protect\phantomsection\label{citation-network-coverage-gaps}
The 2,176 intra-corpus edges spanning 547 connected components provide a
topological skeleton, but three systematic gaps inflate the component
count: (1) cross-source identifier mismatches (DOI vs.~OpenAlex
vs.~arXiv ID), (2) papers whose references are not indexed by any source
API, and (3) open-access preprints whose DOIs differ from their
published versions. Exhaustive DOI-level cross-matching with fuzzy title
matching would condense the graph further.
\end{block}

\begin{block}{Corpus Biases, Citation Dynamics, and Linguistic Framing}
\protect\phantomsection\label{corpus-biases-citation-dynamics-and-linguistic-framing}
Citation counts are subject to Matthew effects and cumulative field-size
biases. Partial-year indexing for the most recent calendar year
undercounts recent publications. The measured 20.36\% CAGR reflects the
dilutive effect of the long temporal span (2005--2026), corresponding to
a 2.8-year publication doubling time; the growth phase from 2010 onward
follows a steeper trajectory with a substantially shorter doubling
interval. Additionally, the retrieved corpus itself suffers from
selection biases inherent to queried databases, including
English-language dominance and the structural over-indexing of preprints
relative to peer-reviewed final versions. Finally, the predominantly
positive hypothesis scores across the board are inflated by two
systematic effects: (1) \textbf{publication bias}, which causes academic
journals to preferentially select positive and confirmatory findings
\citep{sterling1959publication}, and (2) \textbf{linguistic asymmetry}
in scientific writing, where declarative claims are phrased
affirmatively far more often than negatively. These effects jointly
suppress contradicting assertions in the extracted evidence base.
Relative rankings and temporal trajectories are more reliable than
absolute score magnitudes.
\end{block}

\begin{block}{LLM Extraction Fidelity, Domain Drift, and Robustness}
\protect\phantomsection\label{llm-extraction-fidelity-domain-drift-and-robustness}
Zero-shot LLM extraction introduces distinct systematic biases:
over-extraction (the model hallucinating certainty for claims the paper
merely mentions in passing) and direction inversion (misclassifying
opposing evidence as supporting). Recent benchmarking confirms that
state-of-the-art systems often fall short of production-level precision
on tasks requiring exhaustive retrieval and aggregation of directional
claims from long documents \citep{liang2024survey}. Furthermore, because
our corpus extends to 2026, LLM extraction is vulnerable to \emph{domain
drift}---the base models may lack parametric knowledge of the most
recent theoretical developments. As an alternative, fine-tuned models
specifically trained on FEP/AIF abstracts could yield higher precision
than our zero-shot approach, though at a steeper computational setup
cost.

To formally bound these extraction errors and ensure robustness, a
formal validation protocol is required, including a 10\%
manual-annotation ground-truth baseline evaluated via Cohen's
\(\kappa\). Such validation will verify that the automated scoring
pipeline meets minimum inter-rater reliability thresholds
(\(\kappa > 0.70\)) before aggregating the data. The explicit
``irrelevant'' filtering predicate further mitigates over-extraction,
converting what would be a vulnerability of automated reviews into a
calibrated evidential ledger.
\end{block}
\end{frame}

\begin{frame}[fragile,allowframebreaks]{Research Agenda: Four Priority
Next Steps \label{sec:next_steps}}
\protect\phantomsection\label{research-agenda-four-priority-next-steps}
The current prototype establishes a reproducible baseline and surfaces
the field's evidence structure at corpus scale. Four concrete next
steps, ordered by the dependency chain each one unlocks, define the path
from prototype to production-grade living review.

\begin{block}{Next Step 1 --- Expand the Scope of Referenced Data}
\protect\phantomsection\label{next-step-1-expand-the-scope-of-referenced-data}
The present corpus of \(N = 817\) papers is retrieved via keyword
queries against three APIs (Semantic Scholar, OpenAlex, arXiv). Three
expansion axes would materially change the evidence landscape.

\textbf{Additional sources.} PubMed, PsycINFO, and IEEE Xplore each
index Active Inference literature that the current APIs do not reach:
neuroscience clinical trials (PubMed), cognitive-behavioral studies
(PsycINFO), and robotics control architectures (IEEE). For each new
source, the retrieval layer requires only a source-specific connector
implementing the same \breaktt{fetch\_papers(query,\ max\_results)}
interface used by existing adapters. Gray literature---technical
reports, theses, and institutional preprints not yet indexed by major
APIs---represents an additional tier: harvesting from ORCID work records
and institutional repositories would capture practitioner findings that
never appear in indexed venues.

\textbf{Broader query coverage.} The current query set is derived from
the eight hypothesis keywords and their immediate synonyms. Expanding to
a full ontological synonym set (e.g., mapping ``variational inference,''
``surprise minimization,'' and ``Helmholtz machine'' as equivalent
retrieval terms for FEP-related claims) would reduce the retrieval
false-negative rate for papers that use non-canonical vocabulary. A
systematic evaluation of retrieval precision and recall against a
hand-curated gold-standard set of 100 known AIF papers would quantify
the gap.

\textbf{Custom curated bibliographies.} Domain experts can contribute
citation lists directly to the corpus without modifying any code:
placing a \texttt{.bib} or \texttt{.ris} file in
\breaktt{data/custom\_bibliographies/} triggers the deduplication merge
on the next pipeline run. This pathway is the lowest-friction route to
extending scope for researchers who maintain personal reference
libraries.
\end{block}

\begin{block}{Next Step 2 --- Extract and Verify Evidence Supporting
Claims in Each Paper}
\protect\phantomsection\label{next-step-2-extract-and-verify-evidence-supporting-claims-in-each-paper}
The current extraction pipeline operates exclusively on abstracts.
Abstracts contain the claims authors choose to foreground, not
necessarily the claims best supported by the paper's data. Three
mechanisms bridge this gap.

\textbf{Full-text ingestion.} For the subset of papers with open-access
PDFs (approximately 60--70\% of recent AIF preprints on arXiv), Stage 3
can be extended to parse full-text sections---specifically Methods,
Results, and Discussion---using a structured chunking strategy that
splits documents into \textasciitilde512-token segments aligned to
section boundaries. The existing nanopublication schema accommodates a
\texttt{source\_section} field (currently unused) that would record the
provenance of each extracted assertion (abstract vs.~results
vs.~discussion), enabling downstream stratification of evidence by
rhetorical function.

\textbf{Claim-evidence pairing.} The current extraction prompt asks the
LLM to classify a paper's stance toward a hypothesis but does not
require it to quote the specific sentence or data point that justifies
the classification. A revised prompt would require the model to (a)
identify the hypothesis-relevant passage verbatim, (b) classify the
stance, and (c) rate confidence on the basis of whether the passage
reports an empirical measurement, a theoretical derivation, or an
assertion without quantitative support. This three-field extraction ---
\texttt{evidence\_quote}, \texttt{stance}, \texttt{evidence\_type} ---
upgrades the nanopublication from a classification label to a traceable
evidential pointer. For H3 (Markov Blanket Realism), where the 108
contradicting assertions drive a contested score, reviewers could then
inspect the actual quoted passages rather than trusting the LLM
classification in isolation.

\textbf{Human spot-check coverage.} The planned 10\% manual-annotation
baseline would focus on inter-rater agreement; this manual validation
has not yet been performed. Extending spot-checks to verify that the
extracted evidence quote actually appears in the source document (a
verbatim-match check) adds an additional fidelity gate beyond stance
accuracy.
\end{block}

\begin{block}{Next Step 3 --- Tie Hypotheses to Real-World Outcomes}
\protect\phantomsection\label{next-step-3-tie-hypotheses-to-real-world-outcomes}
The eight tracked hypotheses are formulated at the level of theoretical
constructs (e.g., ``the FEP provides a universal account of
self-organizing systems''). Practical applicability requires mapping
from hypothesis support to observable real-world outcomes,
distinguishing which claims are actionable from which remain theoretical
scaffolding.

\textbf{Outcome taxonomy.} Each hypothesis should be annotated with a
set of \emph{outcome indicators}: specific, measurable real-world
results whose observation would constitute evidence for or against the
hypothesis under the closest empirical operationalization. For example:
- H4 (Predictive Coding): outcome indicator = reduction in
prediction-error amplitude as measured by ERP N400 or oscillatory
gamma-band response in human neuroimaging studies. - H5 (Scalability):
outcome indicator = task performance on standard RL benchmarks (Atari,
MuJoCo, ProcGen) at or above the performance of model-free SOTA at
matched computational budgets. - H6 (Clinical Utility): outcome
indicator = statistically significant improvement on standardized
psychiatric assessment scales (PANSS, BDI-II, PTSD Checklist) in at
least one registered clinical trial. - H7 (Morphogenesis): outcome
indicator = quantitative recapitulation of at least one morphogenetic
patterning sequence (e.g., digit formation timecourse, limb bud size
scaling) in a computational model governed by FEP dynamics rather than
reaction-diffusion equations.

For each hypothesis, the extraction pipeline can be extended to tag
assertions whose evidence type is \breaktt{empirical\_measurement} and
whose outcome aligns with these indicators, producing a filtered score
that counts only outcome-linked evidence. This \emph{outcome-filtered
score} sits alongside the current citation-weighted score in the
hypothesis table, providing a direct answer to ``how much of this
support is grounded in real-world observations rather than theoretical
commentary?''

\textbf{Application domain cross-walk.} The subfield classification
(A1--C5) already partitions the corpus by application domain.
Intersecting hypothesis scores with application domain
membership---computing \(\text{score}(H_i, D_j)\) for each hypothesis
\(H_i\) and domain \(D_j\)---would reveal which domains are generating
empirical traction versus theoretical citation counts. H1 (FEP
Universality) likely has high A2 (philosophy) support and lower C1--C5
empirical support; quantifying this split would replace qualitative
description of the ``neutral plurality'' with a decomposed evidence
profile grounded in domain labels already computed by Stage 2.
\end{block}

\begin{block}{Next Step 4 --- Formal Evaluation Rubric for Pipeline
Quality}
\protect\phantomsection\label{next-step-4-formal-evaluation-rubric-for-pipeline-quality}
The current validation is primarily structural: do scripts run, do
outputs exist, do tests pass, does the PDF render? A formal evaluation
rubric answers a different question: \emph{how accurate is the evidence
landscape this system produces?} Four rubric dimensions, together with
their measurement protocols and target thresholds, define what ``good
enough for a published living review'' means.

\begin{table}[htbp]
\centering
\caption{Proposed evaluation rubric for pipeline quality assessment. Each dimension has a measurement protocol, a current baseline, and a target threshold for a production living review. All metrics are computed on a held-out annotation set of 200 randomly sampled assertions.}
\label{tab:eval_rubric}
\begin{tabular}{llll}
\toprule
\textbf{Dimension} & \textbf{Protocol} & \textbf{Current} & \textbf{Target} \\
\midrule
Extraction direction accuracy & Cohen's $\kappa$ (rule reference vs.\ LLM stance) & $\kappa = -0.055$ & $\kappa > 0.80$ \\
Evidence-quote fidelity & Verbatim substring match rate & 0.085 (no quotes stored) & $\geq 90\%$ \\
Corpus recall & Precision/recall vs.\ rule reference & $P = 0.000$, $R = 0.000$ & recall $\geq 0.85$ \\
Outcome grounding rate & Fraction of supporting assertions citing an outcome indicator & pending & $\geq 30\%$ \\
\bottomrule
\end{tabular}
\end{table}

The four rubric dimensions map directly to the four next steps: corpus
recall measures Step 1 progress, evidence-quote fidelity measures Step 2
progress, outcome grounding rate measures Step 3 progress, and
extraction direction accuracy is the targeted baseline to be established
as the other three improve. Reporting all four numbers alongside
hypothesis scores in each pipeline release converts a qualitative
description of limitations into a versioned, trackable quality
scorecard. This transforms the current ``we acknowledge limitations''
posture into an audit trail: readers can see whether the rubric scores
improved between release v1.0 and v2.0, and reviewers can evaluate
pipeline trustworthiness on principled criteria rather than subjective
judgement.
\end{block}
\end{frame}

\begin{frame}[allowframebreaks]{Future Directions: Beyond Tally-Based
Evidence Aggregation}
\protect\phantomsection\label{future-directions-beyond-tally-based-evidence-aggregation}
Beyond the four priority next steps above, the scoring machinery itself
can be upgraded. We identify four directions, ordered by expected
impact.

\begin{block}{Hierarchical Bayesian Hypothesis Scoring}
\protect\phantomsection\label{hierarchical-bayesian-hypothesis-scoring}
The most direct extension replaces the additive tally with a
\textbf{hierarchical Bayesian model} that treats each hypothesis score
as a latent variable inferred from noisy assertion observations. Under
this formulation, each assertion \(a_i\) contributes a likelihood term
\(P(a_i | \theta_H, \sigma)\) parameterized by the hypothesis-level
evidence strength \(\theta_H\) and an observation noise term \(\sigma\)
capturing LLM extraction uncertainty. A hierarchical prior
\(\theta_H \sim \mathcal{N}(\mu_{\text{field}}, \tau^2)\) pools
information across hypotheses, enabling principled shrinkage for
hypotheses with sparse evidence (e.g., H6 Clinical Utility, with only 38
assertions). This framework produces posterior credible intervals rather
than point estimates, providing uncertainty quantification that the
current tally-based scores lack. Temporal dynamics can be modeled
through time-varying parameters \(\theta_H(t)\) using state-space
formulations that re-weight older evidence rather than treating all
cumulative assertions equally.
\end{block}

\begin{block}{Causal Evidence Graphs}
\protect\phantomsection\label{causal-evidence-graphs}
A second-generation knowledge graph would encode not only
assertion-level relationships (paper → supports → hypothesis) but also
\textbf{causal dependencies among hypotheses} themselves. For example,
evidence for predictive coding (H4) often implicitly supports FEP
universality (H1), yet the tally-based approach treats them as
independent. A causal evidence graph---structured as a directed acyclic
graph (DAG) over hypotheses with edge weights learned from co-assertion
patterns---would enable cross-hypothesis evidence propagation using
belief propagation or variational message passing. This is particularly
relevant for the Active Inference literature, where hypotheses are
theoretically nested: FEP universality (H1) logically entails predictive
coding (H4), and Markov blanket realism (H3) is a prerequisite for
certain formulations of H1. Encoding these dependencies would prevent
the double-counting of evidence from papers that support multiple
related hypotheses and enable identification of which specific claims
drive support for downstream hypotheses. The resulting causal structure
itself would be a scientific contribution---a formal map of evidential
dependencies within the field's theoretical architecture.
\end{block}

\begin{block}{Evidential Diversity and Source Weighting}
\protect\phantomsection\label{evidential-diversity-and-source-weighting}
The current formula weights assertions by
\(\log(1 + \text{citations}) \cdot \text{confidence}\), treating all
assertion sources symmetrically. A more nuanced approach would introduce
an \textbf{evidential diversity index} that downweights correlated
evidence from papers sharing authors, institutions, or methodological
approaches. Concretely, assertions could be weighted by the inverse of
their similarity to previously counted assertions, measured via cosine
similarity of paper embeddings. This would address the observation that
H1 (FEP universality) accumulates a large neutral tally partly because
many A2 (philosophy) papers invoke the FEP without independently testing
it---a form of evidential redundancy that inflates the evidence base
without adding independent information. Additionally, assertions could
be stratified by evidence type (empirical, theoretical, review) with
configurable type-specific weights, enabling users to compute evidence
scores that privilege experimental results over theoretical commentary.
\end{block}

\begin{block}{Agentic LLM Extractors and Domain Adaptation}
\protect\phantomsection\label{agentic-llm-extractors-and-domain-adaptation}
Drawing on recent work extending active inference into artificial
reasoning \citep{friston2025active} and proposing AIF as a reward-free
alternative for LLM-based agents \citep{wen2025missing}, replacing
static prompt templates with goal-directed reasoning architectures could
significantly improve confidence calibration. As demonstrated by Friston
et al.~\citep{friston2025active}, ``active reasoning'' enables agents to
perform structure learning---determining which causal rule governs a
situation by seeking observations that explicitly disambiguate competing
hypotheses about world models. Applied to literature extraction,
analogous uncertainty-aware reasoning could treat each paper as a
structured observation to be parsed against hypothesis definitions via
an optimal experimental design rubric---directly operationalizing Next
Step 2's claim-evidence pairing at scale. The framework is
domain-agnostic by design; adaptation to foundation models, quantum
computing, or synthetic biology requires only domain-specific hypothesis
definitions and keyword lists within the A/B/C taxonomy. The broader
convergence between AIF and deep learning demonstrated by AXIOM
\citep{heins2025axiom}---which plans in object-centric
state-spaces---further validates this trajectory. Systematic
cross-referencing with the Energy-Based Model research program
\citep{lecun2006tutorial}---including Helmholtz machines
\citep{dayan1995helmholtz}, contrastive divergence training
\citep{hinton2002training}, and variational autoencoders
\citep{kingma2014auto}---would illuminate shared mathematical structures
currently obscured by disciplinary siloing.
\end{block}
\end{frame}

\begin{frame}[allowframebreaks]{Limitations}
\protect\phantomsection\label{limitations}
Three constraints bound the current findings. First, extraction operates
on abstracts only: full-text methods, results, and supplementary
data---where quantitative effect sizes and experimental controls
live---are not yet parsed. The rubric's evidence-quote fidelity
dimension (Step 4, Table\textasciitilde{}\ref{tab:eval_rubric}) will
quantify exactly how much signal this omission suppresses once a
full-text pilot is run. Second, keyword-based retrieval across three
APIs produces a snapshot with systematic false negatives: papers using
non-canonical terminology, gray literature, and domain-adjacent work
(EBM, Bayesian brain models) are undercounted. The corpus recall metric
provides a principled bound on this gap rather than a vague
acknowledgement of it. Third, the citation-weighted tally treats all
assertion sources symmetrically; the evidential diversity and
outcome-grounding extensions above are the concrete remedies. These are
not general disclaimers but tracked deficits against which the Step 1--4
roadmap makes measurable progress.
\end{frame}

\begin{frame}[allowframebreaks]{Broader Impact}
\protect\phantomsection\label{broader-impact}
Knight et al.~\citep{knight2022fep} identified three capabilities as
goals for the field: ``encompass increased scope of relevant works,''
``integrate multiple forms of annotation and participation,'' and
``facilitate integration of manual and artificial contributions.'' The
four-step research agenda in §\ref{sec:next_steps} operationalizes each
of these directly: Step 1 addresses scope, Step 2 addresses the quality
of extracted contributions, Step 3 addresses empirical grounding, and
Step 4 provides the formal rubric that makes ``integration of manual and
artificial contributions'' verifiable rather than aspirational.

By demonstrating that LLM-driven assertion extraction can produce
scalable, queryable representations of scientific evidence---processing
\(N = 817\) papers spanning approximately two and a half decades
(2005--2026), extracting structured assertions, and evaluating 8 core
hypotheses---this work provides a reusable architecture for realizing
this vision. The corpus window begins in 2005 to capture Energy-Based
Model and variational Bayesian antecedents that predate the Free Energy
Principle label itself; the formal FEP was introduced in 2006
\citep{friston2006free} and reached its core elaboration by 2010
\citep{friston2010free}. The citation network metrics (2,176 edges,
0.33\% density, mean in-degree 2.7) characterize the field's structure,
which has grown at a 20.36\% CAGR while diversifying across 5
application domains.

The limitations of keyword-based retrieval across disjoint academic
repositories mean that any retrieved corpus will contain both false
positives and false negatives. There is no single threshold that
perfectly defines inclusion or exclusion for a dynamic,
interdisciplinary research field. The primary contribution of this work
is therefore not a definitive corpus but an open-source, modularly
updatable, and versioned software package. This tool is built in
reference to custom literature bibliographies that can be iteratively
curated for relevance by the community.

The combination of multi-source retrieval, LLM-based extraction, and
probabilistic knowledge graph construction provides a reusable template
that advances each of these goals. A complementary pathway is emerging
through Retrieval-Augmented Generation (RAG) architectures that ground
LLMs directly in knowledge graphs, reducing hallucination and enabling
real-time, context-aware reasoning over structured evidence
\citep{fan2024survey}. Integrating our nanopublication graph into such a
RAG system would enable natural-language querying of the evidence base,
further lowering the barrier for community engagement. The recent
release of nanopub-js v0.1.0 \citep{kuhn2026nanopubjs}---enabling
browser-based creation, signing, and querying of
nanopublications---lowers the barrier for community-contributed
assertions, bringing the participatory evidence curation envisioned by
Knight et al.~within practical reach. As LLM capabilities improve and
standardized metadata adoption grows, the cost of maintaining such
systems will decrease while their utility increases. By open-sourcing
the pipeline and publishing the schema, we provide both a concrete tool
for the Active Inference community and a modular blueprint that other
fields can adapt and refine.

\textbf{Data and code availability.} The pipeline source code,
configuration, and manuscript templates are available in the project
repository (see \breaktt{metadata.repository} in \texttt{config.yaml} or
the manuscript front matter). Nanopublications are persisted as JSON
Lines (for incremental runs) and RDF/TriG (nanopub.net-compliant); both
can be archived with the code release or on a data repository (e.g.,
Zenodo) for citation and long-term access.

Community recommendations, actionable implications, and open questions
arising from this work are detailed in the
\hyperref[sec:discussion]{Discussion}.
\end{frame}

\end{document}
